\subsection{Two minimal columns and at most two outside elements}

\begin{lemma}[Proved: two-dimensional coefficients for a bounded two-term complex]
\label{CP2:A-two-dimensional-two-term-coefficients}
Let $A$ be a nonzero unital rational algebra of dimension at most two,
let $V$ be any finite-dimensional left $A$-module, and let $D$ be a
complex of finite-dimensional right $A$-modules in degrees zero and one.
If $D\otimes_A^{\mathbf L}V$ has bounded homology, then
\[
 H_j(D\otimes_A^{\mathbf L}V)=0\qquad(j\ge2).
\]
\end{lemma}
\begin{proof}
If $A$ is semisimple, $V$ is projective and the conclusion is immediate.
Otherwise $\dim A=2$.  Choosing a basis $1,a$ shows that $A$ is
$\mathbb Q[t]/(f)$ for a monic quadratic $f$.
In characteristic zero a square-free $f$ gives a product of fields,
so the nonsemisimple case has $f=(t-c)^2$, with $c\in\mathbb Q$.
Consequently $A=\mathbb Q[\epsilon]/(\epsilon^2)$.

The square-zero linear operator $\epsilon$ on $V$ has Jordan blocks
of sizes one and two.  A size-two block is the regular module $A$,
and a size-one block is $k=A/(\epsilon)$.  Thus
\[
 V\cong A^{\oplus a}\oplus k^{\oplus b}.
\]
The regular summands contribute only degrees zero and one.
If $b>0$, boundedness of the whole derived tensor implies boundedness
of $D\otimes_A^{\mathbf L}k$, because this is a direct summand.
Lemma~\ref{CP2:A-local-two-term-derived-amplitude} applies to the
character with nilpotent kernel $(\epsilon)$ and kills its homology
in every degree at least two.  This proves the result for any $V$,
without a one-dimensional coefficient hypothesis.
\end{proof}

\begin{theorem}[Proved: two minimal columns with at most two outside elements]
\label{CP2:A-two-column-two-outside-graph}
Let $M$ be a finite monoid with proper rectangular minimal band
\[
 K=\{e_{ij}:0\le i<r,\ j\in\{0,1\}\},\qquad
 e_{ij}e_{uv}=e_{iv},
\]
where $r\ge1$, and suppose $|M\setminus K|\le2$.
Let $S$ be a finite nonempty left-$N$, right-$M$ biset.
Suppose $BN$ is rationally acyclic and its action on $S$ has the
natural equivariant augmented graph-chain comparison used in
Theorem~\ref{CP2:A-projective-local-row-graph-obstruction}.
If the collage has bounded rational homology, all its rational
degree-two products vanish.

In particular the conclusion applies to $N=R_3$ and every finite
connecting biset.  The number $r$ of minimal rows and the number of
states are unrestricted.
\end{theorem}
\begin{proof}
Write $R=\mathbb QM$, $A=R/\mathbb QK$, $e=e_{00}$, and $e'=e_{01}$.
Collapsing the actual rationally acyclic corner $BN$ gives a rational
homology equivalence to the actual cofiber
\[
 Y=\operatorname{hocofib}(B(N,S,M)\longrightarrow BM).
\]
It therefore induces a rational cohomology-ring isomorphism.
The graph comparison is used only for its natural augmented chains.
With its right-$M$ graph denoted by $\Gamma$, put
\[
 C=[C_1(\Gamma;\mathbb Q)\longrightarrow
       \ker(C_0(\Gamma;\mathbb Q)\to\mathbb Q)]
\]
in degrees one and zero.  The reduced chains of the actual $Y$ are
computed by $(C\otimes_R^{\mathbf L}\mathbb Q)[1]$.

Suppose first that the image graphs $\Gamma e$ and $\Gamma e'$
have a common vertex $v$.  Both retractions fix $v$, and
$ve_{ij}=vee_{ij}=ve_{0j}=v$ for every band element.
For every $m\in M$, the ideal condition gives $em\in K$, so
$vm=vem=v$.  The resulting augmentation section of graph chains
is $M$-linear.  Transport through the natural comparison makes the
actual bridge pullback injective on $H^2(BM;\mathbb Q)$.
Every second class on $Y$ is therefore relative to $BM$.
Such classes come from the actual suspension quotient of the
cofiber and all their products are zero.

In the remaining case the image graphs have disjoint vertex sets
and hence disjoint edge sets.  All band actors of column $j$ have
the same image as $e_{0j}$: the equalities
$e_{ij}e_{0j}=e_{ij}$ and $e_{0j}e_{ij}=e_{0j}$ give both inclusions.
Thus multiplication is a degreewise injective chain map
\[
 J=Ce\otimes_{\mathbb Q}eR\longrightarrow C,
 \qquad \operatorname{im}(J)=C\mathbb QK.
\]
Indeed $eR$ has exactly the two column basis elements $e,e'$;
the two copies of $Ce$ have disjoint supports.  In degree zero
both copies have augmentation zero, as required by the definition
of $C$.  Every term of $J$ is a projective right $R$-module.
The quotient $D=C/J$ is consequently a two-term right $A$-complex.
In degrees at least three, removing $J$ changes no derived tensor
homology, since $J\otimes_R^{\mathbf L}\mathbb Q$ is itself two-term.

The left projective $Re$ has basis $e_{i0}$, and its augmentation
has kernel the row-augmentation module
\[
 0\longrightarrow V\longrightarrow Re\longrightarrow\mathbb Q
 \longrightarrow0,
 \qquad V=\mathbb Q\{e_{i0}-e_{00}:1\le i<r\}.
\]
The band kills $V$, so $V$ is a left $A$-module.
Since $De=0$, derived tensoring this sequence identifies
$D\otimes_R^{\mathbf L}\mathbb Q$ with
$(D\otimes_R^{\mathbf L}V)[1]$.
Multiplication also gives
\[
 \mathbb QK\cong Re\otimes_{\mathbb Q}eR.
\]
This ideal is left projective.  Its tensor with $D$ is acyclic,
because $De=0$.  The resulting left-projective resolution
$0\to\mathbb QK\to R\to A\to0$ therefore yields
$D\otimes_R^{\mathbf L}A\simeq D$, and derived associativity
identifies $D\otimes_R^{\mathbf L}V$ with
$D\otimes_A^{\mathbf L}V$.
The two shifts give
\[
 H_n(Y;\mathbb Q)=H_{n-2}(D\otimes_A^{\mathbf L}V)
 \qquad(n\ge4).
\]
Boundedness of $Y$ implies boundedness of the tensor on the right.
The algebra $A$ has dimension $|M\setminus K|\le2$, so the preceding
lemma applies to the full row module $V$, of any dimension $r-1$.
It gives $H_4(Y;\mathbb Q)=0$, which kills all degree-two products
in this second case as well.

For $N=R_3$, the natural graph comparison is exactly the one in
Corollary~\ref{CP2:A-baseline-right-zero-three-six}.  Its construction
depends only on the opposite $R_3$ action, not on the source order,
and $BR_3$ is contractible.  Thus all the required hypotheses hold.
\end{proof}

\begin{corollary}[Proved: seven remaining three-row order-eight sources]
\label{CP2:A-eight-three-row-two-column-sources}
The order-eight source indices
\[
 501,\ 503,\ 504,\ 505,\ 506,\ 508,\ 509
\]
in the frozen positive minimal-band catalog cannot be the source of
an eighteen-arrow $\mathbb CP^2$ collage opposite $R_3$.
\end{corollary}
\begin{proof}
Each literal table has a $3\times2$ minimal band and exactly two
outside elements, including the identity.  The exact table and
band-coordinate links are independently reconstructed in
\path{converse/scope18/two_column_two_outside_source_scope.json}.
Apply the theorem, since a $\mathbb CP^2$ model has bounded
rational homology.  The assertion uses no finite action census.
\end{proof}
